\begin{helper}[Transport of a contained stored-pair cycle]
Let \(t<\epsilon\le n\) be slots of a certified Tashkinov series.  Assume
slot \(t\) is PE with pair \(\{\gamma_t,\delta_t\}\), every slot strictly
between \(t\) and \(\epsilon\) is RE, and slot \(n\) is PE.  Let
\(\sigma\) be a proper, exactly-one-uncolored coloring stable from the
current coloring at \(n\) on its preceding tree with the current stored
set.  If a full \((\gamma_t,\delta_t)\)-component under the current
coloring at \(t\) is a cycle and all of its vertex-nodes lie in the
preceding tree at \(\epsilon\), then its full component under \(\sigma\)
has exactly the same support and is again a cycle.
\end{helper}
\begin{proof}
The PE constructor at \(t\) in \noderef{TashkinovSeriesStep} identifies
its stored pair with \(\{\delta_t,\gamma_t\}\).  Hence each of the two
colors is either missing on the preceding tree at \(t\), or survives in
\noderef{TashkinovStoredColors} there.  Repeatedly apply the two
inclusions of \noderef{HkuMissingStoredMonotone}: first cross each
coloring step and then each ordered tree extension.  Both colors thus
belong to the missing-or-stored control set at the preceding tree of
\(\epsilon\), and subsequently to the corresponding terminal control
set at \(n\).

Every intermediate slot is RE, whose current-coloring equation in
\noderef{TashkinovSeriesStep} fixes the coloring.  Therefore the current
coloring at \(t\) equals the preceding coloring at \(\epsilon\).  Choose
the cycle witnessing \noderef{TwoColorComponentCycle}.  If an exact
edge-node belongs to its support, each endpoint vertex-node is adjacent
to it in \noderef{TwoColorIncidenceGraph}, hence belongs to the same
full component.  The containment hypothesis puts both endpoints in the
tree at \(\epsilon\), so the edge belongs to \noderef{InternalEdges}.
The internal-edge clause of \noderef{SameColorSeriesPreserved}, applied
from \(\epsilon\) through \(n\), therefore fixes the exact color of every
cycle edge.  Thus the same alternating cycle occurs under the current
coloring at \(n\).

The PE certificate at \(n\) makes that current coloring proper and
exactly-one-uncolored.  At every cycle vertex there is already one
incident cycle edge of each of the two distinct colors.  Properness in
the exact-edge sense of \noderef{IsProperEdgeColoring} excludes every
additional incident edge of either color.  Consequently the displayed
cycle support is the entire two-color component under the current
coloring at \(n\), not merely a subcycle.

The cycle lies in the tree at \(n\) by the ordered-extension clauses of
the series, and both colors lie in its missing-or-stored control set as
shown above.  Forward \noderef{StableColoring} therefore fixes every
cycle-edge identity from the current coloring at \(n\) to \(\sigma\).
For the reverse inclusion, package the two proper,
exactly-one-uncolored colorings and use symmetry from \noderef{sc2}.
Vertexwise missing sets agree on the tree, so the same two colors are in
the reverse control set.  If an edge incident with a cycle vertex had
\(\sigma\)-color \(\gamma_t\) or \(\delta_t\), reverse stability would
give it the same color under the current coloring at \(n\).  Properness
there forces it to be the already fixed cycle edge of that color.
Therefore no \(\sigma\)-two-color path can leave the cycle support, while
the fixed cycle edges keep the support connected.  The two full
components are equal, and the same closed walk supplies the required
\noderef{TwoColorComponentCycle} under \(\sigma\).
\end{proof}
